\documentclass[11pt]{article}

\usepackage[T1]{fontenc}
\usepackage{lmodern}
\usepackage[margin=1.08in]{geometry}
\usepackage{amsmath,amssymb,amsthm,mathtools}
\usepackage{microtype}
\usepackage[backend=biber,style=numeric-comp,sorting=nyt,maxbibnames=99]{biblatex}
\addbibresource{references.bib}
\usepackage[colorlinks=true,linkcolor=blue,citecolor=blue,urlcolor=blue]{hyperref}
\usepackage[nameinlink,capitalise,noabbrev]{cleveref}
\usepackage{enumitem}
\usepackage{booktabs}

\hypersetup{
  pdftitle={Conjugate fibres, certificate rigidity, and irrationality exponents},
  pdfauthor={Christopher D. Long},
  pdfsubject={Irrationality exponents of algebraic exponential fibres}
}
\setlist{topsep=4pt,itemsep=3pt}
\setlength{\emergencystretch}{2em}

\newcommand{\CC}{\mathbb C}
\newcommand{\RR}{\mathbb R}
\newcommand{\QQ}{\mathbb Q}
\newcommand{\ZZ}{\mathbb Z}
\newcommand{\NN}{\mathbb Z_{\ge 0}}
\newcommand{\Qbar}{\overline{\mathbb Q}}
\newcommand{\OO}{\mathcal O}
\newcommand{\eps}{\varepsilon}
\newcommand{\Gm}{\mathbb G_m}
\newcommand{\Comp}{\Sigma_{\xi}(a,b)}
\DeclareMathOperator{\Hom}{Hom}
\DeclareMathOperator{\Aut}{Aut}
\DeclareMathOperator{\Gal}{Gal}
\DeclareMathOperator{\lcm}{lcm}
\DeclareMathOperator{\ord}{ord}
\DeclareMathOperator{\Tr}{Tr}
\DeclareMathOperator{\Nm}{N}
\DeclareMathOperator{\PGL}{PGL}
\DeclareMathOperator{\arctanh}{arctanh}

\newtheorem{theorem}{Theorem}[section]
\newtheorem{proposition}[theorem]{Proposition}
\newtheorem{lemma}[theorem]{Lemma}
\newtheorem{corollary}[theorem]{Corollary}
\newtheorem{conjecture}[theorem]{Conjecture}
\theoremstyle{definition}
\newtheorem{definition}[theorem]{Definition}
\theoremstyle{remark}
\newtheorem{remark}[theorem]{Remark}
\numberwithin{equation}{section}
\crefname{conjecture}{conjecture}{conjectures}
\Crefname{conjecture}{Conjecture}{Conjectures}

\title{Conjugate fibres, certificate rigidity, and irrationality exponents}
\author{Christopher D. Long\\[2pt]\small \texttt{galizur@gmail.com}}
\date{Draft of October 13, 2026}

\begin{document}
\maketitle

\begin{abstract}
Let $\xi\in\RR^\times$ and let $a,b\in\Qbar^\times$ satisfy
$e^{a\xi}=b$.  We first prove two rigidity statements.  Every algebraic
exponential certificate for a fixed $\xi$ has coefficient in one
rational line, and every compatible embedding of the minimal
certificate field is either the identity or a unique inversion
involution.  Thus, if $F=\QQ(a,b)$ and $d=[F:\QQ]$, there are at most two
compatible embeddings.  We prove
\[
  \mu(\xi)\le 2d,
\]
improved to $\mu(\xi)\le d$ when the inversion involution exists.
Equivalently, if $r\in\{1,2\}$ is the number of compatible embeddings,
then
\[
  \mu(\xi)\le \frac{2d}{r}.
\]
The proof combines the bounded-fibre interpolation determinant of
Liu--Jiang--Zhang with simultaneous number-field denominator clearing
and the product formula.  At the compatible places we jointly optimize
the linear rational-approximation cost and the quadratic Taylor-collision
energy.  This removes the usual low/high transverse-index dichotomy and
identifies $2d/r$ as the exact threshold of the resulting aggregate scalar
estimate.

Rigidity reduces optimization over all certificates to power saturation
of any one certificate.  It produces an intrinsic integer
$\kappa_{\exp}(\xi)\ge1$, attained by a sufficiently divisible power,
such that
\[
  \mu(\xi)\le2\kappa_{\exp}(\xi).
\]
We classify $\kappa_{\exp}(\xi)=1$: it is precisely the union of split
rational logarithms and logarithm coordinates of rational points on
quadratic norm-one tori, including the completed quadratic
$\arctanh/\arctan$ family and its torsion boundary.  Beyond complexity
one we obtain, among other consequences, the optimal scalar bound
$\mu(\log\alpha)\le2d_{\mathrm{pow}}(\alpha)$ for positive algebraic
$\alpha$, relative norm-one bounds, and finite bounds for algebraic
multiples of $\pi$ and arguments of algebraic numbers.
\end{abstract}

\section{Introduction}\label{sec:introduction}

For an irrational real number $x$, its \emph{irrationality exponent} is
\[
 \mu(x)=\sup\left\{\nu>0:
  0<\left|x-\frac pq\right|<q^{-\nu}
  \text{ for infinitely many }p\in\ZZ,
  \ q\in\ZZ_{>0}\right\}.
\]
Dirichlet's theorem gives $\mu(x)\ge2$ for every irrational $x$; see,
for example, \cite[Chapter~I]{Cassels1957}.  Roth's theorem gives equality
for every irrational algebraic number \cite{Roth1955}.  For explicit
transcendental constants, proving exponent two is generally much harder.

OpenAI recently proved $\mu(\pi)=2$ using a separated-weight
interpolation determinant \cite{OpenAI2026}.  Two independent subsequent
works moved the interpolation centres between fibres of the exponential
coordinate.  Brown proved the rational-logarithm theorem and supplied a
Lean~4 formalization of that theorem; he also recorded in
\cite[Proposition~5.1]{Brown2026} the bound
$\mu(\log\alpha)\le2[\QQ(\alpha):\QQ]$ for positive real algebraic
$\alpha$.  Liu, Jiang, and Zhang proved
\[
  \mu(\log q)=2
  \qquad(q\in\QQ_{>0}\setminus\{1\})
\]
and established the bounded-fibre moving-centre interpolation theorem
that we use below \cite{LiuJiangZhang2026}.  Their paper explicitly
isolates the next arithmetic obstruction: for an algebraic argument,
clearing denominators produces a nonzero algebraic integer rather than a
nonzero rational integer, so estimates are needed at its conjugate
embeddings.  The first purpose of this note is to carry out that
comparison whenever some conjugate exponential fibres remain compatible.

Litsey recently converted the determinant proof of the irrationality of
Catalan's constant into a finite irrationality measure by controlling mixed
determinants under column perturbations and by retaining the energy cost of
configurations with missing points \cite[Theorem~5.1 and
Proposition~5.10]{Litsey2026}.  The compatible-place estimate below uses a
related continuous-defect principle: rather than separating low and high
transverse indices, it optimizes the approximation and collision penalties
simultaneously.  The argument itself is applied directly to the
row-translation and collision bounds of Liu--Jiang--Zhang.

Eric Naslund publicly pointed out, in posts on X on 9 and 12 October
2026, the exact-two quadratic family
\[
 \frac{1}{\sqrt D}\arctanh(t\sqrt D),
 \qquad t\in\QQ,
\]
with the appropriate real branches and torsion boundary; the later post
linked a public Claude discussion \cite{Naslund2026}.  For non-torsion
members of this family, the moving-fibre interpolation supplied by Brown
and Liu--Jiang--Zhang is needed in addition to the fixed-fibre argument
for $\pi$.  The new structural results here are the converse rigidity
statement, the intrinsic power-saturation theorem, and the higher-degree
bounds beyond that quadratic exact-two locus.

Fix $\xi\in\RR^\times$ and $a,b\in\Qbar^\times$ with $e^{a\xi}=b$.
Let
\[
 F=\QQ(a,b),\qquad d=[F:\QQ],\qquad
 \Sigma_F=\Hom_{\QQ}(F,\CC),
\]
and define
\begin{equation}\label{eq:compatible-set}
 \Comp=
 \left\{\sigma\in\Sigma_F:
   e^{\sigma(a)\xi}=\sigma(b)\right\}.
\end{equation}
The distinguished inclusion belongs to $\Comp$, so
$1\le r:=|\Comp|\le d$.

\begin{theorem}[Rigid conjugate-fibre bound]\label{thm:main}
Let $\xi\in\RR^\times$ and $a,b\in\Qbar^\times$ satisfy $e^{a\xi}=b$,
and put $F=\QQ(a,b)$ and $d=[F:\QQ]$.
\begin{enumerate}[label=\textup{(\roman*)}]
\item One always has
\[
 \boxed{\mu(\xi)\le2d.}
\]
\item If $F$ has an embedding $\tau:F\hookrightarrow\CC$ satisfying
\[
 \tau(a)=-a,
 \qquad
 \tau(b)=b^{-1},
\]
then $\tau$ is a unique involution and
\[
 \boxed{\mu(\xi)\le d.}
\]
\end{enumerate}
Equivalently, if $r$ denotes the number of compatible embeddings, then
$r\in\{1,2\}$ and
\[
 \boxed{\mu(\xi)\le\frac{2d}{r}.}
\]
More precisely, for every real $\nu>2d/r$ there is
$Q=Q(a,b,\xi,\nu)$ such that
\[
  \left|\xi-\frac pq\right|\ge q^{-\nu}
  \qquad(p\in\ZZ,\ q\in\ZZ_{>0},\ q\ge Q).
\]
In particular, $\mu(\xi)=2$ when $d=1$, or when $d=2$ and the inversion
involution exists.
\end{theorem}

The exact-two assertions use Dirichlet's theorem and
Hermite--Lindemann.  If $\xi$ were nonzero algebraic, then $a\xi$ would
be nonzero algebraic, whereas $e^{a\xi}=b$ is algebraic, a contradiction;
see \cite[Chapter~1]{Baker1975}.

The second purpose of the note is to show that the scalar compatibility
in \cref{thm:main} is rigid.  If $e^{c\xi}$ is algebraic for another
nonzero algebraic coefficient $c$, then $c/a\in\QQ^\times$.  Moreover,
every compatible embedding of $F$ sends
\[
 (a,b)\longmapsto(a,b)
 \quad\text{or}\quad
 (a,b)\longmapsto(-a,b^{-1}).
\]
Thus $r\le2$.  This turns the optimized scalar theory into an
integer-valued one: after replacing $b$ by sufficiently divisible
powers, the ratios $d/r$ decrease to a positive integer
$\kappa_{\exp}(\xi)$, and
\[
 \mu(\xi)\le2\kappa_{\exp}(\xi).
\]
The case $\kappa_{\exp}(\xi)=1$ is classified by one-dimensional tori:
the split torus gives rational logarithms, while a quadratic norm-one
torus gives the $\arctanh/\arctan$ family and its period boundary.

The paper is organized as follows.  \Cref{sec:rigidity} proves the
rigidity theorems.  Sections~\ref{sec:interpolation}--\ref{sec:proof}
prove \cref{thm:main}.  \Cref{sec:complexity} constructs the intrinsic
integer complexity, and \cref{sec:classification} classifies complexity
one.  \Cref{sec:applications} gives applications, and
\cref{sec:conjectures} records conjectural strengthenings.

\section{Rigidity of algebraic exponential certificates}\label{sec:rigidity}

We use the Gelfond--Schneider theorem in the following standard form:
if $\alpha\in\Qbar\setminus\{0,1\}$ and
$\beta\in\Qbar\setminus\QQ$, then every value of $\alpha^\beta$ is
transcendental; see \cite[Chapter~1]{Baker1975}.

\begin{theorem}[Certificate-line rigidity]\label{thm:certificate-line}
Let $\xi\in\RR^\times$, and suppose
\[
 e^{a\xi}=b,
 \qquad
 e^{c\xi}=u,
 \qquad
 a,b,c,u\in\Qbar^\times.
\]
Then
\[
 \boxed{\frac ca\in\QQ^\times.}
\]
Thus the coefficients of all algebraic exponential certificates for a
fixed $\xi$ lie on one rational line.
\end{theorem}

\begin{proof}
Put $q=c/a\in\Qbar^\times$.  Suppose first that $b\ne1$.  Since
$a\xi$ is a logarithm of $b$,
\[
 u=e^{c\xi}=e^{q a\xi}
\]
is one value of $b^q$.  If $q$ were algebraic irrational,
Gelfond--Schneider would make $u$ transcendental.  Hence $q\in\QQ$.

Suppose now that $b=1$.  Since $a\ne0$ and $\xi\ne0$, there is a
nonzero integer $n$ such that
\[
 a\xi=2\pi i n.
\]
Then
\[
 u=e^{2\pi i n q}
\]
is the value of $(-1)^{2nq}$ obtained from the logarithm
$\log(-1)=\pi i$.  If $q$ were algebraic irrational,
Gelfond--Schneider would again make $u$ transcendental.  Thus
$q\in\QQ^\times$ in all cases.
\end{proof}

\begin{theorem}[Compatible-fibre rigidity]\label{thm:compatible-rigidity}
Let $e^{a\xi}=b$ and $F=\QQ(a,b)$ as above.  If
$\sigma:F\hookrightarrow\CC$ is compatible, then
\[
 \boxed{
  (\sigma(a),\sigma(b))=(a,b)
  \quad\text{or}\quad
  (\sigma(a),\sigma(b))=(-a,b^{-1}).}
\]
The first possibility is the distinguished inclusion.  The second, if it
occurs, is a unique involution in $\Aut(F/\QQ)$.  In particular,
\[
 \boxed{r\in\{1,2\}.}
\]
\end{theorem}

\begin{proof}
Put $q=\sigma(a)/a$.  If $b\ne1$, compatibility gives
\[
 \sigma(b)=e^{\sigma(a)\xi}=e^{q a\xi},
\]
so the same Gelfond--Schneider argument as in
\cref{thm:certificate-line} gives $q\in\QQ^\times$.  If $b=1$, write
$a\xi=2\pi i n$ and $\sigma(a)\xi=2\pi i m$ with nonzero integers
$m,n$; then $q=m/n\in\QQ^\times$.

Extend $\sigma$ to an automorphism $g$ of a normal closure of $F$.
Such an extension exists by the standard extension theorem for field
embeddings; see, for example, \cite[Chapter~I]{Neukirch1999}.  If $N$
is the finite order of $g$, then
\[
 a=g^N(a)=q^N a.
\]
Hence $q^N=1$.  Since $q\in\QQ^\times$, one has $q=1$ or $q=-1$.

If $q=1$, compatibility gives $\sigma(a)=a$ and $\sigma(b)=b$, so
$\sigma$ fixes the generators of $F$ and is the distinguished
inclusion.  If $q=-1$, then
\[
 \sigma(a)=-a,
 \qquad
 \sigma(b)=e^{-a\xi}=b^{-1}.
\]
These values on the generators determine at most one embedding.  Its
image is contained in $F$, so it is an automorphism of $F$; applying it
twice fixes both generators, hence it is an involution.
\end{proof}

\begin{remark}[Why a second compatible fibre occurs]
\label{rem:complex-conjugation}
Because $\xi$ is real, complex conjugation is always a compatible
embedding:
\[
 e^{\overline a\,\xi}=\overline b.
\]
By \cref{thm:compatible-rigidity}, either
$\overline a=a$ and $\overline b=b$, in which case the distinguished
certificate field is real, or
$\overline a=-a$ and $\overline b=b^{-1}$, in which case $a$ is purely
imaginary and $|b|=1$.  Thus nonreal certificates obtain their second
compatible fibre from complex conjugation.  For real certificates, a
second fibre can only come from a real inversion involution; in the
complexity-one case this is necessarily a real quadratic involution.
\end{remark}

\begin{corollary}\label{cor:integer-ratio}
For every certificate, the ratio $d/r$ is a positive integer.  More
precisely, if $r=1$ then $d/r=d$, while if $r=2$ and $\tau$ is the
inversion involution, then
\[
 \frac dr=\frac d2=[F^\tau:\QQ].
\]
\end{corollary}

\begin{proof}
Only the case $r=2$ requires comment.  Then $\tau$ is a nontrivial
involution of $F$, so $[F:F^\tau]=2$.
\end{proof}

\section{The interpolation input}\label{sec:interpolation}

We recall the precise form of the geometric input that is needed.  The following is \cite[Theorem~9]{LiuJiangZhang2026}.  Fix
positive rational weights
\[
 W=(w_0,w_1,\ldots,w_m),\qquad
 V=(v_0,w_1/\theta,\ldots,w_m/\theta),
 \qquad 0<\theta<1.
\]
Let $\mathcal P_W(H)$ be the span of the monomials $Y^hX^\alpha$ with
\[
 w_0h+\sum_{i=1}^m w_i\alpha_i\le H,
\]
and let $\mathcal J_V(H)$ be the quotient of
$\CC[[t,u_1,\ldots,u_m]]$ retaining precisely the monomials
$t^su^\beta$ with
\[
 v_0s+\frac1\theta\sum_{i=1}^m w_i\beta_i<H.
\]

\begin{theorem}[Bounded-fibre moving-centre interpolation]\label{thm:MI}
Fix $m,K\ge1$, $1\le\varkappa\le K$, positive rational $w_0,v_0$, and
$0<\theta<1$ such that
\begin{equation}\label{eq:MI-volume}
 K\left(\frac{w_0}{v_0}\right)\theta^m<1,
 \qquad \varkappa\theta^m<1.
\end{equation}
There are successive lower thresholds for the positive rational weights
$w_1,\ldots,w_m$, independent of the interpolation centres, with the
following property.  Let
\[
 (y_j,c_{j1},\ldots,c_{jm})\in\CC^\times\times\CC^m,
 \qquad 0\le j<K,
\]
assume that at most $\varkappa$ indices have a given value of $y_j$, and
assume that
\begin{equation}\label{eq:fibre-separation}
 j\ne\ell,\ y_j=y_\ell
 \quad\Longrightarrow\quad
 c_{ji}\ne c_{\ell i}\quad(1\le i\le m).
\end{equation}
Then, for every sufficiently large admissible degree $H$, the map
\[
 \mathcal P_W(H)\longrightarrow
 \bigoplus_{j=0}^{K-1}\mathcal J_V(H)
\]
obtained by substituting
\[
 Y=y_j(1+t),\qquad
 X_i=c_{ji}+u_i+\log(1+t)
\]
is onto.  The divisibility modulus for $H$ depends only on the weights;
the eventual threshold may depend on the centres.
\end{theorem}

The proof in \cite{LiuJiangZhang2026} uses a local intersection-length
estimate, separated weights, a curve inequality, an ordinary blow-up,
and Serre vanishing.  Brown independently proved a general-fibre version
under the global condition $K\theta^m<1$
\cite[Theorem~2.1]{Brown2026}.  We quote the bounded-fibre formulation
of Liu--Jiang--Zhang because its second volume constraint depends only on
the maximal fibre multiplicity $\varkappa$, which is sharper when many
multiplicative fibres are distinct and still covers torsion.  We shall
use the determinant translation and collision estimates from
Sections~4--6 of Liu--Jiang--Zhang.  Those estimates are uniform over the
selected full-row minor and its columns.

\subsection{Parameters}

The optimized collision-energy estimate permits a particularly simple
parameter choice.

\begin{lemma}[Soft-gap parameters]\label{lem:parameters}
Let $d,r$ be positive integers and let $\nu>2d/r$.  There are rational
numbers $\chi,\theta,B,C$ satisfying
\begin{equation}\label{eq:param-soft}
 0<\chi<1,
 \qquad \theta=\chi^2,
 \qquad C=\chi^{-1},
 \qquad \theta<B<\chi,
\end{equation}
and
\begin{equation}\label{eq:param-gap}
 g:=r\nu(B-\theta)-d(1-\theta)>0.
\end{equation}
In particular,
\begin{equation}\label{eq:param-consequences}
 C\theta=\chi<1,
 \qquad CB=\frac B\chi<1,
 \qquad \frac{C\theta}{B}=\frac\chi B>1,
 \qquad B^2<\theta.
\end{equation}
\end{lemma}

\begin{proof}
Since $r\nu>2d$, one may choose a rational $\chi\in(0,1)$ sufficiently
close to $1$ that
\[
 r\nu\chi>d(1+\chi).
\]
This is equivalent to
\[
 \chi^2+\frac{d}{r\nu}(1-\chi^2)<\chi.
\]
Choose a rational number $B$ strictly between the two sides, and put
$\theta=\chi^2$ and $C=\chi^{-1}$.  Then
\[
 r\nu(B-\theta)>d(1-\theta),
\]
which is \eqref{eq:param-gap}; all statements in
\eqref{eq:param-consequences} are immediate.
\end{proof}

For a positive integer $m$, set
\begin{equation}\label{eq:dimension-parameters}
 K=\lfloor C^m\rfloor,
 \qquad w_0=B^{-m},
 \qquad v_0=2K\theta^m w_0.
\end{equation}
Then
\begin{equation}\label{eq:parameter-limits}
 K\left(\frac{w_0}{v_0}\right)\theta^m=\frac12,
 \qquad K\theta^m\longrightarrow0,
 \qquad \frac K{w_0}\longrightarrow0,
 \qquad v_0\longrightarrow\infty,
 \qquad \frac m{v_0}\longrightarrow0.
\end{equation}
Indeed, for all sufficiently large $m$,
\[
 K\theta^m\le(C\theta)^m=\chi^m,
 \qquad
 \frac K{w_0}\le(CB)^m=\left(\frac B\chi\right)^m,
\]
while $K\ge C^m/2$ gives
\[
 v_0\ge\left(\frac{C\theta}{B}\right)^m
      =\left(\frac\chi B\right)^m.
\]

\section{The number-field determinant}\label{sec:determinant}

Fix the data of \cref{thm:main}, and suppose for contradiction that
there is a real $\nu>2d/r$ for which
\begin{equation}\label{eq:good-approximations}
 \left|\xi-\frac pq\right|\le q^{-\nu}
\end{equation}
holds for arbitrarily large positive integers $q$.  Replacing $\nu$ by
a slightly smaller rational number still larger than $2d/r$, we may
assume that $\nu$ is rational.

Choose the parameters of \cref{lem:parameters}.  Fix a rational
truncation parameter
\begin{equation}\label{eq:Phi}
 \Phi>\frac2\theta.
\end{equation}
We shall first take $m$ large, then choose $m$ approximations
\[
 x_i=\frac{p_i}{q_i},\qquad
 \left|\xi-x_i\right|\le q_i^{-\nu},
 \qquad w_i=\lceil\log q_i\rceil,
 \qquad 1\le i\le m,
\]
successively.  The denominators are taken large enough that $x_i\ne0$,
$|x_i|\le|\xi|+1$, and the weights meet the successive thresholds in
\cref{thm:MI}.  This is possible because those thresholds are
independent of the centres.

Use the centres
\begin{equation}\label{eq:centres}
 y_j=b^j,\qquad c_{ji}=ja x_i,
 \qquad 0\le j<K,\quad1\le i\le m.
\end{equation}
We apply \cref{thm:MI} with $\varkappa=K$.  The second inequality in
\eqref{eq:MI-volume} holds for large $m$ by $K\theta^m\to0$.  If
$y_j=y_\ell$ and $j\ne\ell$, then
\[
 c_{ji}-c_{\ell i}=(j-\ell)a x_i\ne0
\]
for every $i$, so \eqref{eq:fibre-separation} holds.  Thus the same
argument covers roots of unity and non-torsion values of $b$.

Set
\begin{equation}\label{eq:truncations}
 T_i=\left\lceil\frac{\Phi w_i}{v_0}\right\rceil,
 \qquad
 G_i(t)=\sum_{1\le k<T_i}\frac{(-1)^{k+1}}{k}t^k,
 \qquad
 L_i=\lcm(1,\ldots,T_i-1),
\end{equation}
where an empty sum and empty least common multiple are interpreted as
$0$ and $1$.  Because $v_0T_i>w_i/\theta$, replacing $\log(1+t)$ by
$G_i(t)$ is an invertible filtered change on every retained jet packet
\cite[Section~4]{LiuJiangZhang2026}.

For an admissible degree $H$, use the monomials
\[
 P=Y^hX^\alpha,
 \qquad w_0h+\sum_iw_i\alpha_i\le H,
\]
as columns, and use the triples
\[
 \varrho=(j,s,\beta),\qquad
 0\le j<K,
 \qquad v_0s+\frac1\theta\sum_iw_i\beta_i<H,
\]
as rows.  The matrix entry is
\begin{equation}\label{eq:matrix-entry}
 [t^su^\beta]\,
 b^{jh}(1+t)^h
 \prod_{i=1}^m(ja x_i+G_i(t)+u_i)^{\alpha_i}.
\end{equation}
For every sufficiently large admissible $H$, interpolation supplies a
nonzero square minor $\Delta_H\in F^\times$ containing every row.  Let
$M=M_H$ be its size, and put
\begin{equation}\label{eq:bbar}
 \bar\beta_H=\frac1{MH}
 \sum_{\varrho=(j,s,\beta)}\sum_iw_i\beta_i.
\end{equation}
The row volume and its first transverse moment will both be useful.

\begin{lemma}[Row volume and first moment]\label{lem:row-moment}
As $H\to\infty$ through admissible degrees,
\begin{equation}\label{eq:row-asymptotic}
 M\ge
 \frac{K\theta^mH^{m+1}}
 {(m+1)!v_0\prod_iw_i},
 \qquad
 M\sim
 \frac{K\theta^mH^{m+1}}
 {(m+1)!v_0\prod_iw_i},
\end{equation}
and
\begin{equation}\label{eq:row-first-moment}
 \bar\beta_H=\frac{m}{m+2}\theta+O(H^{-1}),
 \qquad 0\le\bar\beta_H\le\theta.
\end{equation}
The implied constant depends only on the fixed weights.
\end{lemma}

\begin{proof}
Put
\[
 q_0=v_0,
 \qquad q_i=\frac{w_i}{\theta}\quad(1\le i\le m),
\]
and let
\[
 \mathcal R_H=
 \left\{n=(n_0,\ldots,n_m)\in\NN^{m+1}:
       \sum_{i=0}^m q_i n_i<H\right\}.
\]
The row set is the product of $\{0,\ldots,K-1\}$ with
$\mathcal R_H$.  The unit cubes based at the points of $\mathcal R_H$
cover the simplex
\[
 \left\{x\in\RR_{\ge0}^{m+1}:\sum_iq_ix_i<H\right\},
\]
which gives the lower bound in \eqref{eq:row-asymptotic}.  Comparing the
same cubes with the two parallel simplices obtained by shifting the
boundary by the fixed amount $\sum_iq_i$ gives
\[
 \#\mathcal R_H
 =\frac{H^{m+1}}{(m+1)!\prod_{i=0}^m q_i}+O(H^m).
\]
The same comparison, now with the linear functions $q_i x_i$, gives
\[
 \sum_{n\in\mathcal R_H}q_i n_i
 =\frac{H^{m+2}}{(m+2)!\prod_{j=0}^m q_j}+O(H^{m+1})
 \qquad(0\le i\le m).
\]
Since
\[
 \sum_iw_i\beta_i
 =\theta\sum_{i=1}^m q_i n_i,
\]
dividing the corresponding first-moment asymptotic by
$H\#\mathcal R_H$ gives
$m\theta/(m+2)+O(H^{-1})$.  Finally, every row satisfies
$w\beta<\theta H$, which gives the stated upper bound.
\end{proof}

\begin{remark}[Limiting row distribution]\label{rem:row-beta}
Under uniform counting of rows, the normalized transverse weight
$w\beta/(\theta H)$ converges to the $\operatorname{Beta}(m,2)$
distribution.  Equation \eqref{eq:row-first-moment} is its first moment.
\end{remark}

\subsection{Algebraic-integral clearing and the norm inequality}

Choose an integer $D_0\ge1$ such that
\[
 D_0a,\ D_0b\in\OO_F.
\]
Put
\[
 \mathfrak D_H=\prod_{i=1}^mL_i^{\lfloor H/w_i\rfloor}.
\]
Scale the column $(h,\alpha)$ and the row $(j,s,\beta)$ by
\begin{equation}\label{eq:scales}
 C_{h,\alpha}=D_0^{(K-1)h}\prod_i(D_0q_i)^{\alpha_i},
 \qquad
 S_{j,s,\beta}=\mathfrak D_H\prod_i(D_0q_i)^{-\beta_i}.
\end{equation}

\begin{lemma}[Norm lower bound]\label{lem:norm}
Let $\Lambda=4\log2$ and $w_*=\min_iw_i$.  Define
\begin{equation}\label{eq:Ear}
 E_{\mathrm{ar}}=
 \frac{\Lambda\Phi m}{v_0}
 +\Lambda\sum_{i=1}^m\frac1{w_i}
 +\frac{\theta+\log D_0}{w_*}
 +\frac{(K-1)\log D_0}{w_0}.
\end{equation}
Then
\begin{equation}\label{eq:norm-inequality}
 0\le
 \sum_{\sigma\in\Sigma_F}
 \frac{\log|\sigma(\Delta_H)|}{MH}
 +d\bigl(1-\bar\beta_H+E_{\mathrm{ar}}\bigr).
\end{equation}
\end{lemma}

\begin{proof}
The entry obtained from \eqref{eq:matrix-entry} after the row and column
scalings is zero unless $\beta\le\alpha$.  Otherwise it equals
\begin{align*}
 &\mathfrak D_H(D_0b)^{jh}D_0^{(K-1-j)h}
 \binom\alpha\beta[t^s](1+t)^h\\
 &\hspace{25mm}\cdot
 \prod_{i=1}^m
 \bigl(j(D_0a)p_i+D_0q_iG_i(t)\bigr)^{\alpha_i-\beta_i}.
\end{align*}
The denominator contributed by the $i$th truncated logarithm divides
$L_i^{\alpha_i-\beta_i}$, and
$\alpha_i-\beta_i\le\lfloor H/w_i\rfloor$.  Hence every scaled entry
belongs to $\OO_F$.  The scaled determinant
\[
 \widetilde\Delta_H
 =\Delta_H\prod_{\text{columns}}C_{h,\alpha}
             \prod_{\text{rows}}S_{j,s,\beta}
\]
is therefore a nonzero algebraic integer.  Consequently
\[
 1\le|\Nm_{F/\QQ}(\widetilde\Delta_H)|
 =\prod_{\sigma\in\Sigma_F}|\sigma(\widetilde\Delta_H)|;
\]
we count complex embeddings separately, as in the usual product formula
\cite[Chapter~I]{Neukirch1999}.

It remains to bound the positive rational scaling factor.  As in
\cite[Section~4]{LiuJiangZhang2026},
\[
 \frac{\log\mathfrak D_H}{H}
 \le\frac{\Lambda\Phi m}{v_0}
    +\Lambda\sum_i\frac1{w_i}.
\]
Every selected column satisfies
\[
 \sum_i\alpha_i\log q_i\le H,
 \qquad |\alpha|\le\frac H{w_*},
 \qquad h\le\frac H{w_0}.
\]
Moreover, since $0\le w_i-\log q_i<1$,
\[
 \sum_{\text{rows}}\sum_i\beta_i\log q_i
 \ge MH\bar\beta_H-\frac{MH\theta}{w_*}.
\]
The negative row contribution involving $\log D_0$ may be discarded.
Thus the logarithm of the total scaling factor, divided by $MH$, is at
most $1-\bar\beta_H+E_{\mathrm{ar}}$.  Taking logarithms of the norm
inequality proves \eqref{eq:norm-inequality}.
\end{proof}

\section{Archimedean determinant estimates}\label{sec:analytic}

Set
\[
 C_a=\max_{\sigma\in\Sigma_F}\max\{1,|\sigma(a)|\},
 \qquad
 C_b=\max_{\sigma\in\Sigma_F}\max\{1,|\sigma(b)|\},
 \qquad X=|\xi|+1.
\]
All constants in this section are uniform over the finite set
$\Sigma_F$.  The verification below is included to make explicit that
neither the row translation nor the collision argument requires real
centres or distinct multiplicative fibres.  Reality enters only through
the common rational approximants to $\xi$; repeated fibres are handled
solely by the interpolation theorem that supplies the full-row minor.

\subsection{Compatible embeddings}

For $\sigma\in\Comp$, define
\[
 z_j(t)=j\sigma(a)\xi+\log(1+t),
 \qquad
 \epsilon_{ji}=j\sigma(a)(x_i-\xi),
 \qquad
 \tau_i(t)=G_i(t)-\log(1+t).
\]
Compatibility gives the exact identities
\begin{equation}\label{eq:compatible-identities}
 e^{z_j(t)}=\sigma(b)^j(1+t),
 \qquad
 j\sigma(a)x_i+G_i(t)+u_i
 =z_j(t)+u_i+\epsilon_{ji}+\tau_i(t).
\end{equation}
For a selected column $P=Y^hX^\alpha$ and a multi-index $\lambda$, let
\[
 f_{\lambda,P}(z)
 =[u^\lambda]P(e^z,z+u_1,\ldots,z+u_m)
 =\binom\alpha\lambda e^{hz}z^{|\alpha|-|\lambda|},
\]
with the convention that this is zero unless $\lambda\le\alpha$.
Expanding \eqref{eq:compatible-identities} gives the same exact row
translation as in \cite[Section~5.1]{LiuJiangZhang2026}:
\begin{align}\label{eq:row-translation}
 \operatorname{row}_{j,s,\beta}^{\sigma}
 ={}&\sum_{\substack{\lambda\ge\beta\\w\lambda\le H}}
 \binom\lambda\beta
 \sum_{0\le e\le\lambda-\beta}
 \binom{\lambda-\beta}{e}
 \epsilon_j^{\lambda-\beta-e}\\
 &\quad\cdot\sum_{k=0}^s
 [t^k]\prod_i\tau_i(t)^{e_i}
 \bigl([t^{s-k}]f_{\lambda,P}(z_j(t))\bigr)_P.
 \nonumber
\end{align}
Here and below $w\lambda=\sum_iw_i\lambda_i$.

Since
\[
 |\epsilon_{ji}|
 \le KC_aq_i^{-\nu}
 \le KC_ae^{\nu}e^{-\nu w_i},
\]
and a nonzero tail coefficient in \eqref{eq:row-translation} uses
$we<H/\Phi$, every scalar multiplying a translated test row is bounded
by
\begin{equation}\label{eq:translation-scalar}
 \exp\{-\nu w(\lambda-\beta)+HE_{\mathrm{tr}}\},
\end{equation}
where one may take
\begin{equation}\label{eq:Etr}
 E_{\mathrm{tr}}
 =\frac\nu\Phi+\frac{\log2}{v_0}
 +\frac{\log4+\log(2KC_a)+\nu}{w_*}.
\end{equation}

Choose
\[
 R=100K(1+C_a|\xi|).
\]
On $|t|\le1/2$, every $z_j(t)$ lies in $|z|<R/2$.  We use the
Taylor-collision argument of \cite[Section~5.2]{LiuJiangZhang2026}; see
also \cite{Laurent2000} for its interpolation-determinant form.  Put
\begin{equation}\label{eq:Ehol}
 E_{\mathrm{hol}}
 =\frac R{w_0}+\frac{\log2}{v_0}
  +\frac{\log(2R)}{w_*},
 \qquad c=\frac{\log2}{4}.
\end{equation}
Let
\[
 S=\sum_iw_i.
\]
The following elementary estimate is the soft replacement for the usual
low/high transverse-index dichotomy.

\begin{lemma}[Collision-energy optimization]\label{lem:collision-energy}
Let $(n_\lambda)_{\lambda\in\NN^m}$ be finitely supported nonnegative
real numbers with $\sum_\lambda n_\lambda=M$.  For every $U>0$,
\begin{equation}\label{eq:soft-energy}
 c\sum_\lambda n_\lambda^2
 +\nu\sum_\lambda w\lambda\,n_\lambda
 \ge
 \nu UHM-
 \frac{\nu^2}{4c}
 \sum_\lambda(UH-w\lambda)_+^2.
\end{equation}
Moreover,
\begin{equation}\label{eq:lattice-second-moment}
 \sum_{\lambda\in\NN^m}(UH-w\lambda)_+^2
 \le
 \frac{2(UH+S)^{m+2}}{(m+2)!\prod_iw_i}.
\end{equation}
\end{lemma}

\begin{proof}
For every $e\ge0$ and $n\ge0$, completion of the square gives
\[
 cn^2+\nu en
 \ge
 \nu UHn-\frac{\nu^2}{4c}(UH-e)_+^2.
\]
Summing with $e=w\lambda$ proves \eqref{eq:soft-energy}.  If
$x\in\lambda+[0,1]^m$, then
\[
 UH+S-wx\ge UH-w\lambda.
\]
Integration over the unit cubes therefore gives
\[
 \sum_\lambda(UH-w\lambda)_+^2
 \le
 \int_{\RR_{\ge0}^m}(UH+S-wx)_+^2\,dx.
\]
After the change of variables $y_i=w_ix_i$, the last integral equals the
right side of \eqref{eq:lattice-second-moment}.
\end{proof}

\begin{lemma}[Optimized compatible-place estimate]\label{lem:compatible}
There is a nonnegative quantity $\delta_H\to0$, independent of
$\sigma\in\Comp$ and of the selected full-row minor, such that for every
$U>0$,
\begin{align}\label{eq:compatible-bound-general}
 \frac{\log|\sigma(\Delta_H)|}{MH}
 \le{}& E_{\mathrm{tr}}+E_{\mathrm{hol}}+\delta_H
 -\nu(U-\bar\beta_H)\\
 &+\frac{\nu^2v_0}{2cK\theta^m(m+2)}
       \left(U+\frac SH\right)^{m+2}.
 \nonumber
\end{align}
Define
\begin{equation}\label{eq:Um-alpha}
 U_m=
 \left(\frac{2cK\theta^m}{\nu v_0}\right)^{1/(m+1)},
 \qquad
 \alpha_m=\frac{m+1}{m+2}U_m.
\end{equation}
Then, for all sufficiently large $H$,
\begin{equation}\label{eq:compatible-bound}
 \frac{\log|\sigma(\Delta_H)|}{MH}
 \le
 E_{\mathrm{tr}}+E_{\mathrm{hol}}+\delta_H
 +\frac{\nu S}{H}
 -\nu(\alpha_m-\bar\beta_H).
\end{equation}
For the parameters in \eqref{eq:dimension-parameters},
\begin{equation}\label{eq:Um-limit}
 U_m=\left(\frac{cB^m}{\nu}\right)^{1/(m+1)},
 \qquad \alpha_m\longrightarrow B.
\end{equation}
\end{lemma}

\begin{proof}
The row expansion \eqref{eq:row-translation} has only polynomially many
terms in $H$ for each row.  Choose one translated row from each original
row and let $n_\lambda$ be the number chosen with transverse index
$\lambda$.  Expanding the entire functions $f_{\lambda,P}$ in Taylor
series in $z/R$, any nonzero determinant contribution uses distinct
Taylor degrees within each fixed $\lambda$-group.  Hence
\begin{equation}\label{eq:collision-term}
 |\det T|
 \le\exp\left\{-c\sum_\lambda n_\lambda^2
 +MH(E_{\mathrm{hol}}+o(1))\right\}.
\end{equation}
To see explicitly why no reality or distinct-fibre hypothesis is hidden
here, write
\[
 f_{\lambda,P}(z)=\sum_{d\ge0}c_{\lambda,d,P}(z/R)^d.
\]
The coefficient vector $(c_{\lambda,d,P})_P$ is independent of the
centre index $j$ and of the row order $s$; all dependence on the complex
point $z_j(t)$ occurs in a scalar coefficient.  Hence two selected rows
with the same pair $(\lambda,d)$ are proportional in the corresponding
determinant term, which is the sole source of the collision saving.
Cauchy's estimate is taken on the complex disk $|z|\le R$, so the
argument is indifferent to whether the points $j\sigma(a)\xi$ are real
or imaginary.  It also does not require distinct multiplicative fibres.
Indeed, if $\sigma(b)^j=\sigma(b)^\ell$ with $j\ne\ell$, then
\[
 z_j(0)-z_\ell(0)=(j-\ell)\sigma(a)\xi\ne0,
\]
although this additive distinctness is not used in the determinant
cancellation.  This is the collision estimate of
\cite[Section~5.2]{LiuJiangZhang2026}; its proof is uniform in the
columns and depends on the centres only through the radius $R$.

Multiplying the row scalars in \eqref{eq:translation-scalar} and using
\[
 \sum_{\text{rows}}w\beta=MH\bar\beta_H
\]
shows that the chosen determinant contribution is at most
\begin{align*}
 \exp\{&MH(E_{\mathrm{tr}}+E_{\mathrm{hol}}+o(1))
       +\nu MH\bar\beta_H\\
      &\hspace{24mm}
       -\nu\sum_\lambda w\lambda\,n_\lambda
       -c\sum_\lambda n_\lambda^2\}.
\end{align*}
Apply \cref{lem:collision-energy}, then use the lower bound for $M$ in
\eqref{eq:row-asymptotic}.  This gives
\eqref{eq:compatible-bound-general}.  The total number of multilinear
expansion terms is polynomial in $H$ to the power $M$, so its logarithm
is $o(MH)$ and is absorbed into $\delta_H$.

For large $H$, take $U=U_m-S/H>0$.  The defining equation for $U_m$
gives
\[
 \frac{\nu^2v_0}{2cK\theta^m(m+2)}U_m^{m+2}
 =\frac{\nu U_m}{m+2}.
\]
Substitution proves \eqref{eq:compatible-bound}.  Finally,
$v_0=2K\theta^mw_0$ and $w_0=B^{-m}$ give
\eqref{eq:Um-limit}.
\end{proof}

\begin{corollary}[Full-Pl\"ucker analytic bound]\label{cor:plucker-analytic}
Let $\mathbf P_H^\sigma$ be the vector of all maximal minors of the
conjugate interpolation matrix at $\sigma\in\Comp$.  For every fixed
$1\le p\le\infty$ there is a nonnegative quantity
$\delta_{H,p}\to0$, uniform in $\sigma\in\Comp$, such that
\begin{equation}\label{eq:plucker-bound}
 \frac1{MH}\log\|\mathbf P_H^\sigma\|_{\ell^p}
 \le
 E_{\mathrm{tr}}+E_{\mathrm{hol}}+\delta_{H,p}
 +\frac{\nu S}{H}
 -\nu(\alpha_m-\bar\beta_H).
\end{equation}
\end{corollary}

\begin{proof}
The preceding proof is uniform in the selected maximal minor.  If
$N_H=\dim\mathcal P_W(H)$, then the number of maximal minors is at most
$2^{N_H}$.  Here $N_H=O(H^{m+1})$, while $MH\asymp H^{m+2}$, so passing
from the largest coordinate to any fixed $\ell^p$-norm costs
$o(MH)$ in the logarithm.
\end{proof}

\begin{remark}[Water-filling profile and stability]\label{rem:water-filling}
The proof retains an exact stability remainder.  For $w\lambda<UH$,
\begin{align*}
 &cn_\lambda^2+\nu w\lambda\,n_\lambda-\nu UHn_\lambda
 +\frac{\nu^2}{4c}(UH-w\lambda)^2\\
 &\qquad=
 c\left(n_\lambda-\frac{\nu}{2c}(UH-w\lambda)\right)^2.
\end{align*}
Thus a contribution close to the optimized bound must have occupancy
close to the truncated affine profile
\[
 n_\lambda\approx\frac{\nu}{2c}(UH-w\lambda)_+.
\]
This identifies precisely what additional incidence or cancellation
information would be required to improve the aggregate estimate.
\end{remark}

\begin{remark}[Robustness of the collision constant]
The precise value $c=\log2/4$ is irrelevant to the limiting threshold.
If the collision estimate in auxiliary dimension $m$ has coefficient
$c_m>0$ and $\log c_m=o(m)$, then
\[
 U_m=B\exp\left(\frac{\log(c_m/(\nu B))}{m+1}\right)
 \longrightarrow B.
\]
Thus any positive collision saving whose deterioration is subexponential
in $m$ leads to the same scalar threshold.
\end{remark}

\begin{remark}[Continuum compatibility of the incidence condition]
At the continuum level, the bare condition $\lambda\ge\beta$ is not
expected to create a leading-order gain.  After summing over the row
order $s$, the scaled row density on the $\theta$-simplex is proportional
to $\theta-\sum_i x_i$, while the water-filling density on the
$U$-simplex is proportional to $U-\sum_i y_i$.  The radial dilation
$y=(U/\theta)x$ carries one profile to a constant multiple of the other
and satisfies $y_i\ge x_i$ when $U>\theta$.  Any further improvement
must therefore use the detailed column incidence, determinant
cancellation, or a non-scalar arithmetic input.
\end{remark}

\subsection{The remaining embeddings}

At an incompatible embedding no approximation gain is available, but
none is needed.  A direct Cauchy estimate gives a normalized cost that
can be made arbitrarily small by the separated parameter choices.

\begin{lemma}[Crude estimate at every embedding]\label{lem:crude}
For every $\sigma\in\Sigma_F$,
\begin{equation}\label{eq:crude-bound}
 \frac{\log|\sigma(\Delta_H)|}{MH}
 \le E_{\mathrm{bad}}+\frac{\log M}{2H},
\end{equation}
where
\begin{equation}\label{eq:Ebad}
 E_{\mathrm{bad}}
 =\frac{K\log C_b+\log(3/2)}{w_0}
 +\frac{\log\bigl(2(1+KC_aX+\log2)\bigr)}{w_*}
 +\frac{\log2}{v_0}.
\end{equation}
\end{lemma}

\begin{proof}
The conjugate of \eqref{eq:matrix-entry} is zero unless
$\beta\le\alpha$.  Otherwise it is
\[
 \sigma(b)^{jh}\binom\alpha\beta[t^s](1+t)^h
 \prod_i(j\sigma(a)x_i+G_i(t))^{\alpha_i-\beta_i}.
\]
On $|t|=1/2$ we have $|G_i(t)|\le\log2$, while
$|x_i|\le X$, $j<K$, $s\le H/v_0$, $h\le H/w_0$, and
$|\alpha|\le H/w_*$.  Cauchy's estimate and
$\binom\alpha\beta\le2^{|\alpha|}$ therefore bound each entry by
$\exp(HE_{\mathrm{bad}})$.  Hadamard's inequality contributes at most
$M^{M/2}$, which gives \eqref{eq:crude-bound}.
\end{proof}

\subsection{Making the errors small}

The choices are made in the following order.  First choose $\Phi$ so
large that the contribution $r\nu/\Phi$ is smaller than any prescribed
positive tolerance.  Next let $m$ tend to infinity.  By
\eqref{eq:parameter-limits} and \eqref{eq:Um-limit},
\[
 \frac{m}{v_0}\to0,
 \qquad \frac K{w_0}\to0,
 \qquad \frac1{v_0}\to0,
 \qquad \alpha_m\to B.
\]
The constants $D_0,C_a,C_b,\xi,d,r$ are fixed.  Hence $m$ may be chosen
so that all terms in \eqref{eq:Ear}, \eqref{eq:Ehol}, and
\eqref{eq:Ebad} not involving $w_i$ are as small as desired and, in
addition,
\begin{equation}\label{eq:soft-gap-finite}
 r\nu(\alpha_m-\theta)-d(1-\theta)>\frac g2.
\end{equation}
Finally, because the interpolation thresholds are centre-independent,
the approximations may be selected successively with all $w_i$ so large
that
\[
 \sum_i\frac1{w_i},
 \qquad
 \frac{1+\log K}{w_*}
\]
are as small as desired.  Once all parameters and centres are fixed,
$H$ tends to infinity through admissible degrees; then
$\delta_H\to0$, $S/H\to0$, and $(\log M)/H\to0$.

\section{Proof of the conjugate-fibre theorem}\label{sec:proof}

Apply \cref{lem:compatible} at the $r$ embeddings in $\Comp$ and
\cref{lem:crude} at the other $d-r$ embeddings.  Substitution in
\eqref{eq:norm-inequality} gives
\begin{equation}\label{eq:final-master}
 0\le
 d(1-\bar\beta_H)-r\nu(\alpha_m-\bar\beta_H)
 +\mathcal E_H,
\end{equation}
where $\mathcal E_H$ is nonnegative.  More explicitly, it is bounded
by the corresponding linear combination of $E_{\mathrm{ar}}$,
$E_{\mathrm{tr}}$, $E_{\mathrm{hol}}$, $E_{\mathrm{bad}}$,
$\delta_H$, $\nu S/H$, and $(\log M)/H$.  The fixed-parameter terms
can be made arbitrarily small by the preceding choices, and the remaining
terms tend to zero with $H$.

Since $r\nu>d$ and $0\le\bar\beta_H\le\theta$,
\begin{align*}
 d(1-\bar\beta_H)-r\nu(\alpha_m-\bar\beta_H)
 &\le d(1-\theta)-r\nu(\alpha_m-\theta)\\
 &<-\frac g2
\end{align*}
by \eqref{eq:soft-gap-finite}.  Taking the fixed errors sufficiently
small and then $H$ sufficiently large contradicts
\eqref{eq:final-master}.

Therefore \eqref{eq:good-approximations} cannot hold for arbitrarily
large denominators when $\nu>2d/r$.  This proves the asserted eventual
lower bound and hence $\mu(\xi)\le2d/r$.  As noted in the introduction,
$\xi$ is transcendental.  When $d=r$---which by rigidity means either
$d=r=1$ or $d=r=2$---Dirichlet's theorem supplies the opposite inequality
$\mu(\xi)\ge2$.  This completes the proof of \cref{thm:main}.

\begin{proposition}[Finite-dimensional barrier of the aggregate estimate]
\label{prop:soft-barrier}
Fix $m\ge1$, $\nu\ge2$, and parameters $B,\theta$ with
$0<B<1$ and $B^2<\theta<1$, and put $c=\log2/4$.  Define
\[
 U_m=\left(\frac{cB^m}{\nu}\right)^{1/(m+1)},
 \qquad
 \alpha_m=\frac{m+1}{m+2}U_m,
 \qquad
 \beta_m=\frac{m}{m+2}\theta.
\]
Then
\begin{equation}\label{eq:finite-barrier}
 2\alpha_m<1+\beta_m.
\end{equation}
Consequently, if $r\nu\le2d$, the limiting leading term supplied by the
optimized aggregate estimate satisfies
\begin{equation}\label{eq:finite-barrier-master}
 d(1-\beta_m)-r\nu(\alpha_m-\beta_m)\ge0.
\end{equation}
\end{proposition}

\begin{proof}
Since $c=\log2/4<\nu$, one has
\[
 U_m<B^{m/(m+1)}<\theta^{m/(2m+2)}.
\]
Set $x=\theta^{1/(2m+2)}\in(0,1)$.  Then
\[
 \alpha_m<\frac{m+1}{m+2}x^m,
 \qquad
 \beta_m=\frac{m}{m+2}x^{2m+2}.
\]
The function
\[
 f(x)=mx^{2m+2}-2(m+1)x^m+m+2
\]
satisfies $f(1)=0$ and
\[
 f'(x)=2m(m+1)x^{m-1}(x^{m+2}-1)\le0
 \qquad(0<x\le1).
\]
Hence $f(x)\ge0$, which proves \eqref{eq:finite-barrier}.  If
$\alpha_m\le\beta_m$, then the left side of
\eqref{eq:finite-barrier-master} is positive.  If
$\alpha_m>\beta_m$, then $r\nu\le2d$ and
\eqref{eq:finite-barrier} give
\[
 d(1-\beta_m)-r\nu(\alpha_m-\beta_m)
 \ge d(1+\beta_m-2\alpha_m)\ge0.
\]
\end{proof}

\begin{remark}[Why the scalar threshold is $2d/r$]
By \cref{lem:row-moment}, $\bar\beta_H\to\beta_m$.  Thus
\cref{prop:soft-barrier} shows that the fully optimized combination of
the linear approximation cost and quadratic collision energy cannot,
for any fixed auxiliary dimension, force a contradiction at or below
$2d/r$.  Conversely, \cref{lem:parameters} produces admissible
parameters for every strict inequality $\nu>2d/r$.  In the limiting
parameter geometry, $B$ may approach $\sqrt\theta$, and
\[
 \frac{1-\theta}{\sqrt\theta-\theta}
 =1+\frac1{\sqrt\theta}\longrightarrow2
 \qquad(\theta\uparrow1).
\]
Thus the factor $2$ comes from the interpolation geometry
$B^2<\theta$, while the factor $d/r$ comes from the scalar
product-formula comparison across conjugate places.  Improvements below
$2d/r$ must use information absent from the aggregate occupancy bound,
such as detailed incidence, cancellation among minors, finite-place
structure, or full-Pl\"ucker arithmetic.
\end{remark}

\begin{remark}[Effectivity]
The theorem is qualitative.  As in \cite{LiuJiangZhang2026}, after the
centres have been selected the analytic contradiction has an explicit
large-degree threshold, but the centre-dependent Serre-vanishing
threshold in the interpolation theorem is not made effective.
\end{remark}


\section{Power saturation and intrinsic complexity}\label{sec:complexity}

Fix one certificate $e^{a\xi}=b$.  For $N\ge1$, put
\begin{equation}\label{eq:power-data-new}
 F_N=\QQ(a,b^N),
 \qquad d_N=[F_N:\QQ],
\end{equation}
and let
\begin{equation}\label{eq:rN-new}
 r_N=\#\left\{\sigma:F_N\hookrightarrow\CC:
 e^{N\sigma(a)\xi}=\sigma(b^N)\right\}.
\end{equation}
Define
\[
 \kappa_N=\frac{d_N}{r_N}.
\]
By \cref{thm:compatible-rigidity,cor:integer-ratio}, each $r_N$ is one
or two and each $\kappa_N$ is a positive integer.

\begin{lemma}[Divisibility monotonicity]\label{lem:divisibility}
If $N\mid M$, then
\[
 \boxed{\kappa_M\le\kappa_N.}
\]
\end{lemma}

\begin{proof}
Since $N\mid M$, one has $F_M\subseteq F_N$.  Every compatible
embedding of $F_N$ restricts to a compatible embedding of $F_M$, because
the relation for $N$ can be raised to the power $M/N$.  Each embedding
of $F_M$ has at most $[F_N:F_M]$ extensions to $F_N$.  Hence
\[
 r_N\le [F_N:F_M]r_M.
\]
Together with $d_N=[F_N:F_M]d_M$, this gives
$d_N/r_N\ge d_M/r_M$.
\end{proof}

For an arbitrary certificate $e^{c\xi}=u$, write
\[
 \kappa(c,u;\xi)=
 \frac{[\QQ(c,u):\QQ]}
 {\#\{\sigma:\QQ(c,u)\hookrightarrow\CC:
 e^{\sigma(c)\xi}=\sigma(u)\}}.
\]

\begin{lemma}[Rational roots do not improve the ratio]
\label{lem:rational-roots}
Let $e^{c\xi}=u$ be any algebraic exponential certificate.  Then there
is a positive integer $M$ such that
\[
 \boxed{\kappa(c,u;\xi)\ge\kappa_M.}
\]
\end{lemma}

\begin{proof}
By \cref{thm:certificate-line}, $c=qa$ for some
$q\in\QQ^\times$.  Replacing $(c,u)$ by $(-c,u^{-1})$ if necessary
does not change its field or compatible-embedding count, so write
$q=m/n$ with coprime positive integers $m,n$.  Then
\[
 u^n=b^m.
\]
Put
\[
 F'=\QQ(a,u),
 \qquad
 F_m=\QQ(a,b^m)\subseteq F',
\]
and let $r'$ be the number of compatible embeddings of $F'$ for the
certificate $e^{c\xi}=u$.  Every such embedding restricts to a
compatible embedding of $F_m$: raising
$e^{(m/n)\sigma(a)\xi}=\sigma(u)$ to the $n$th power gives
$e^{m\sigma(a)\xi}=\sigma(b^m)$.  Therefore
\[
 r'\le[F':F_m]r_m.
\]
Since $[F':\QQ]=[F':F_m]d_m$, one obtains
\[
 \kappa(c,u;\xi)=\frac{[F':\QQ]}{r'}
 \ge\frac{d_m}{r_m}=\kappa_m.
\]
\end{proof}

\begin{definition}[Exponential certificate complexity]
\label{def:kappa-exp}
For a real nonzero $\xi$ admitting an algebraic exponential certificate,
set
\[
 \kappa_{\exp}(\xi)=
 \inf_{e^{c\xi}=u}\kappa(c,u;\xi),
\]
where the infimum runs over all $c,u\in\Qbar^\times$.
\end{definition}

\begin{theorem}[Intrinsic power saturation]\label{thm:intrinsic-kappa}
For any fixed certificate $e^{a\xi}=b$,
\[
 \boxed{
 \kappa_{\exp}(\xi)
 =\min_{N\ge1}\kappa_N
 =\lim_{j\to\infty}\kappa_{j!}
 \in\ZZ_{>0}.}
\]
The sequence $\kappa_{j!}$ is eventually constant.  In particular,
the optimum is attained and is independent of the initial certificate.
Moreover,
\[
 \boxed{\mu(\xi)\le2\kappa_{\exp}(\xi).}
\]
\end{theorem}

\begin{proof}
Powers of the fixed certificate are among all certificates, so
$\kappa_{\exp}(\xi)\le\inf_N\kappa_N$.  Conversely,
\cref{lem:rational-roots} shows that every certificate has ratio at
least one of the $\kappa_N$.  Hence equality holds.

By \cref{lem:divisibility}, the positive integers $\kappa_{j!}$ form a
nonincreasing sequence and are therefore eventually constant.  For
every fixed $N$, one has $N\mid j!$ for all sufficiently large $j$, so
the eventual value is at most $\kappa_N$.  It is therefore the global
minimum.  Finally, apply \cref{thm:main} to a power attaining this
minimum.
\end{proof}

If one defines the optimized compatible-fibre density by
\[
 \rho_{\exp}(\xi)=
 \sup_{e^{c\xi}=u}
 \frac{r(c,u;\xi)}{[\QQ(c,u):\QQ]},
\]
then \cref{thm:intrinsic-kappa} says
\begin{equation}\label{eq:density-reciprocal}
 \boxed{\rho_{\exp}(\xi)=\frac1{\kappa_{\exp}(\xi)}.}
\end{equation}
Thus the only possible optimized densities are
$1,1/2,1/3,\ldots$.

\begin{corollary}[Torsion-defect criterion]\label{cor:torsion-new}
Let $L$ be a normal closure of $\QQ(a,b)$.  Suppose that for every
embedding $\tau:L\hookrightarrow\CC$ the defect
\[
 \delta_\tau=\frac{e^{\tau(a)\xi}}{\tau(b)}
\]
is a root of unity.  Then
\[
 \boxed{\kappa_{\exp}(\xi)=1,
 \qquad \mu(\xi)=2.}
\]
\end{corollary}

\begin{proof}
Choose $N$ divisible by the orders of the finitely many roots of unity
$\delta_\tau$.  Then
\[
 e^{N\tau(a)\xi}=\tau(b^N)
\]
for every $\tau$.  Every embedding of $F_N=\QQ(a,b^N)$ extends to an
embedding of $L$, so every embedding of $F_N$ is compatible.  Hence
$\kappa_N=1$ and \cref{thm:intrinsic-kappa} applies.
\end{proof}

\begin{remark}[Relation lattices]
Suppose $e^{a_j\xi}=b_j$ for finitely many algebraic certificates.
By \cref{thm:certificate-line}, all $a_j$ are rational multiples of
one coefficient.  Any integer combination of the certificates is
therefore another rational rescaling, and
\cref{lem:rational-roots} shows that it cannot improve on power
saturation.  Thus no separate relation-lattice optimization is needed.
\end{remark}

\section{Exact scalar complexity one and one-dimensional tori}
\label{sec:classification}

The point of the next theorem is a converse: Gelfond--Schneider
rigidity shows that, after rational scaling, branch completion, and the
torsion boundary, no other exact-two family can be reached by the scalar
compatible-fibre method.

\begin{theorem}[Classification of $\kappa_{\exp}=1$]
\label{thm:kappa-one-classification}
Let $\xi\in\RR^\times$ admit an algebraic exponential certificate.  The
following are equivalent.
\begin{enumerate}[label=\textup{(\roman*)}]
\item $\kappa_{\exp}(\xi)=1$.
\item Some powered certificate has all of its embeddings compatible.
\item One of the following alternatives holds.
  \begin{enumerate}[label=\textup{(\alph*)}]
  \item There are $\lambda_0\in\QQ^\times$ and
  $\alpha_0\in\QQ_{>0}\setminus\{1\}$ such that
  \[
    \xi=\lambda_0\log\alpha_0.
  \]
  \item There are a quadratic field $E=\QQ(\omega)$ with
  $\omega^2\in\QQ^\times$, its nontrivial automorphism
  $\tau(\omega)=-\omega$, and an element $B\in E^\times$ satisfying
  \[
    \tau(B)=B^{-1},
    \qquad
    e^{2\omega\xi}=B.
  \]
  \end{enumerate}
\end{enumerate}
In alternative \textup{(iii)(b)}, if $B\ne-1$ then
\[
 t=\frac{B-1}{\omega(B+1)}\in\QQ
\]
and, for some $k\in\ZZ$ and a compatible branch of $\arctanh$,
\begin{equation}\label{eq:artanh-parametrization}
 \boxed{
 \xi=\frac1\omega\arctanh(t\omega)
      +\frac{\pi i k}{\omega}.}
\end{equation}
If $B=-1$, then
\begin{equation}\label{eq:torsion-boundary}
 \boxed{
 \xi=\frac{(2k+1)\pi i}{2\omega}}
 \qquad(k\in\ZZ).
\end{equation}
Every number in the theorem has irrationality exponent two.
\end{theorem}

\begin{proof}
The equivalence of \textup{(i)} and \textup{(ii)} follows from the
attainment statement in \cref{thm:intrinsic-kappa}.  Assume
\textup{(ii)}, and choose a powered certificate
\[
 e^{A\xi}=B,
 \qquad
 F=\QQ(A,B),
\]
with every embedding compatible.  By
\cref{thm:compatible-rigidity}, $[F:\QQ]\le2$.

If $F=\QQ$, then $A,B\in\QQ^\times$.  Since $A$ and $\xi$ are real,
$B=e^{A\xi}$ is positive, and $B\ne1$.  Thus
\[
 \xi=A^{-1}\log B,
\]
which is alternative \textup{(iii)(a)}.

Suppose $[F:\QQ]=2$, and let $\tau$ be its nontrivial automorphism.
Compatibility and \cref{thm:compatible-rigidity} give
\[
 \tau(A)=-A,
 \qquad
 \tau(B)=B^{-1}.
\]
Put $\omega=A/2$.  Then $\omega\notin\QQ$,
$F=\QQ(\omega)$, and $\omega^2\in\QQ^\times$.  This is alternative
\textup{(iii)(b)}.

If $B\ne-1$, define $t$ by the displayed formula.  Direct calculation
using $\tau(\omega)=-\omega$ and $\tau(B)=B^{-1}$ gives
$\tau(t)=t$, hence $t\in\QQ$.  Solving for $B$ yields
\[
 B=\frac{1+t\omega}{1-t\omega}.
\]
Since $e^{2\omega\xi}=B$, the logarithms of the two sides differ by
$2\pi i k$, which gives \eqref{eq:artanh-parametrization}.  If $B=-1$,
then $2\omega\xi=(2k+1)\pi i$, giving
\eqref{eq:torsion-boundary}.

Conversely, alternative \textup{(iii)(a)} gives the rational
certificate $e^{\xi/\lambda_0}=\alpha_0$.  In alternative \textup{(iii)(b)}, the
identity and $\tau$ are both compatible with the certificate
$e^{2\omega\xi}=B$.  Hence $\kappa_{\exp}(\xi)=1$ in either case.
The exponent-two conclusion now follows from
\cref{thm:intrinsic-kappa} and Dirichlet's theorem.
\end{proof}

\begin{remark}[Torus interpretation]
Alternative \textup{(iii)(a)} is the split torus $\Gm$.  Alternative
\textup{(iii)(b)} is the quadratic norm-one torus
$\operatorname{Res}^{1}_{E/\QQ}\Gm$; the parameter
$B=(1+t\omega)/(1-t\omega)$ is its Cayley parametrization away from
$B=-1$.  Thus scalar complexity one is exactly the locus of nonzero real
logarithm coordinates of $\QQ$-rational points on one-dimensional
$\QQ$-tori.  For general background on algebraic tori, see
\cite[Chapter~12]{Milne2017}.
\end{remark}

\begin{corollary}[Quadratic artanh family]
\label{cor:quadratic-artanh}
Let $D\in\QQ^\times$, let $t\in\QQ$ with $1-t^2D\ne0$, and choose a branch for which
\[
 \xi=\frac1{\sqrt D}\arctanh(t\sqrt D)
      +\frac{\pi i k}{\sqrt D}
\]
is real and nonzero.  Then
\[
 \boxed{\mu(\xi)=2.}
\]
The torsion boundary
\[
 \xi=\frac{(2k+1)\pi i}{2\sqrt D}
\]
is included as well whenever it is real and nonzero.
\end{corollary}

\begin{proof}
If $\sqrt D\in\QQ$, the value is a rational multiple of a logarithm of
a positive rational number.  Otherwise apply
\cref{thm:kappa-one-classification} with
$\omega=\sqrt D$ and
$B=(1+t\sqrt D)/(1-t\sqrt D)$.
\end{proof}

For $D>0$, the principal real branch has $k=0$ and
$|t|\sqrt D<1$.  Taking $D=1$ and
$t=(\alpha_0-1)/(\alpha_0+1)$ shows that, up to the rational factor $2$, the
family contains every $\log\alpha_0$ with
$\alpha_0\in\QQ_{>0}\setminus\{1\}$.  For $D=-s<0$,
\[
 \frac1{\sqrt D}\arctanh(t\sqrt D)
 =\frac1{\sqrt s}\arctan(t\sqrt s),
\]
so it contains rational arctangents.  Its period boundary contains the
square-root rescalings of $\pi$.

\begin{corollary}[Concrete exact-two values]\label{cor:concrete-values}
For every integer $k\ge2$,
\[
 \mu\!\left(\arctan\frac1k\right)=2,
 \qquad
 \mu\!\left(\frac\pi{\sqrt3}\right)=2,
\]
and
\[
 \boxed{
 \mu\!\left(
  \sqrt{2k-1}\,
  \arctan\frac{\sqrt{2k-1}}{k-1}
 \right)=2.}
\]
\end{corollary}

\begin{proof}
The first assertion is the case $D=-1$, $t=1/k$.  For the second, take
$D=-3$, $t=1$ and multiply the resulting number
$\pi/(3\sqrt3)$ by $3$.  For the last, put
$s=2k-1$ and $t=1/(k-1)$ in the negative-$D$ form above, then multiply
by the rational number $s$.  Irrationality exponent is invariant under
nonzero rational scaling.
\end{proof}

\begin{corollary}[Quadratic norm-one elements]
\label{cor:quadratic-norm-one}
Let $E=\QQ(\sqrt D)$ be real quadratic, and let
$u\in E^\times$ satisfy
\[
 u>0,
 \qquad u\ne1,
 \qquad \Nm_{E/\QQ}(u)=1.
\]
Then
\[
 \boxed{
 \mu\left(\frac{\log u}{\sqrt D}\right)=2.}
\]
If $\Nm_{E/\QQ}(u)=-1$, the same conclusion follows after squaring the
certificate.
\end{corollary}

\begin{proof}
The nontrivial automorphism sends $\sqrt D$ to $-\sqrt D$ and $u$ to
$u^{-1}$.  Apply \cref{thm:kappa-one-classification}; for norm $-1$,
use $u^2$.
\end{proof}

\section{Applications beyond complexity one}\label{sec:applications}

\subsection{Positive algebraic logarithms}

Brown independently obtained the degree bound
$\mu(\log\alpha)\le2[\QQ(\alpha):\QQ]$ for positive real algebraic
$\alpha$ \cite[Proposition~5.1]{Brown2026}.  The theorem below replaces
the degree by the generally smaller power degree and proves that this is
the best possible bound within the entire scalar certificate
architecture.  Liu--Jiang--Zhang already obtain exact exponent two when
a positive power of $\alpha$ is rational
\cite[Corollary~7]{LiuJiangZhang2026}.

For $\alpha\in\Qbar_{>0}\setminus\{1\}$, define
\begin{equation}\label{eq:power-degree-new}
 d_{\mathrm{pow}}(\alpha)=
 \min_{N\ge1}[\QQ(\alpha^N):\QQ].
\end{equation}

\begin{theorem}[Exact scalar complexity of an algebraic logarithm]
\label{thm:alglog-complexity}
For every positive algebraic $\alpha\ne1$,
\[
 \boxed{
 \kappa_{\exp}(\log\alpha)=d_{\mathrm{pow}}(\alpha)}
\]
and consequently
\[
 \boxed{
 \mu(\log\alpha)\le2d_{\mathrm{pow}}(\alpha).}
\]
If some positive power of $\alpha$ is rational, then
$\mu(\log\alpha)=2$.
\end{theorem}

\begin{proof}
Use the initial certificate $e^{\log\alpha}=\alpha$.  For its $N$th
power, the coefficient $N$ is rational and
$F_N=\QQ(\alpha^N)$.  A compatible embedding must satisfy
\[
 \sigma(\alpha^N)=e^{N\log\alpha}=\alpha^N,
\]
so it fixes the generator of $F_N$ and is the identity.  Thus $r_N=1$
and $\kappa_N=[\QQ(\alpha^N):\QQ]$.  Apply
\cref{thm:intrinsic-kappa}.  The rational-power case also follows
directly from \cite[Corollary~7]{LiuJiangZhang2026}.
\end{proof}

\begin{remark}[Algebraic rescaling asymmetry]
The norm-one theorem gives
\[
 \mu\!\left(\frac{\log(3+2\sqrt2)}{\sqrt2}\right)=2,
\]
whereas the scalar complexity of $\log(1+\sqrt2)$ is $2$, so the present
method gives only
\[
 \mu(\log(1+\sqrt2))\le4.
\]
Thus the scalar theory is not invariant under multiplication by a
nonrational algebraic number.  This sharp asymmetry is one motivation
for the full-Pl\"ucker conjectural mechanism discussed in
\cref{sec:conjectures}.
\end{remark}

The power degree has a useful Galois description.

\begin{proposition}[Conjugates modulo torsion]\label{prop:power-degree-galois}
Let $L$ be a normal closure of $\QQ(\alpha)$, put
$G=\Gal(L/\QQ)$, and define
\[
 H_\alpha=
 \left\{g\in G:
 \frac{g(\alpha)}{\alpha}\text{ is a root of unity}\right\}.
\]
Then $H_\alpha$ is a subgroup of $G$ and
\[
 \boxed{d_{\mathrm{pow}}(\alpha)=[G:H_\alpha].}
\]
\end{proposition}

\begin{proof}
The root-of-unity condition is preserved under products and inverses,
so $H_\alpha$ is a subgroup.  Choose $N$ divisible by the orders of all
root-of-unity ratios $g(\alpha)/\alpha$ with $g\in H_\alpha$.  Then
$H_\alpha$ fixes $\alpha^N$.  Conversely, if $g(\alpha^N)=\alpha^N$,
then $g(\alpha)/\alpha$ is an $N$th root of unity, so
$g\in H_\alpha$.  Thus
$H_\alpha=\Gal(L/\QQ(\alpha^N))$, and
$[\QQ(\alpha^N):\QQ]=[G:H_\alpha]$.  This value is minimal over all
powers.
\end{proof}

\subsection{Relative norm-one fibres}

\begin{theorem}[Relative norm-one bound]\label{thm:relative-norm-one}
Let $E$ be a number field, let $\omega,t\in\Qbar$ satisfy
\[
 E=\QQ(\omega^2,t),
 \qquad
 \omega\notin E,
\]
and assume $1-t\omega\ne0$.  Put
\[
 u=\frac{1+t\omega}{1-t\omega}.
\]
Suppose
\[
 \xi=\frac1\omega\arctanh(t\omega)
\]
is real and nonzero for a chosen branch.  Then, with
$F=\QQ(\omega,u)$,
\[
 \boxed{
 \mu(\xi)\le[F:\QQ]\le2[E:\QQ].}
\]
\end{theorem}

\begin{proof}
The extension $E(\omega)/E$ is quadratic with involution
$\tau(\omega)=-\omega$.  Since $t\in E$,
\[
 \tau(u)=u^{-1}.
\]
The relation $e^{2\omega\xi}=u$ therefore has the identity and $\tau$
as compatible embeddings of its minimal certificate field $F$.  Hence
$r=2$, and \cref{thm:main} gives
$\mu(\xi)\le[F:\QQ]$.  Finally
$F\subseteq E(\omega)$.
\end{proof}

The rational quadratic-artanh family is the case $E=\QQ$.  For example,
with $\omega=\sqrt[4]{2}$ and rational $t$ in the real domain,
\[
 \mu\left(
 \frac1{\sqrt[4]{2}}\arctanh(t\sqrt[4]{2})
 \right)\le4.
\]

\subsection{\texorpdfstring{Algebraic multiples of $\pi$}{Algebraic multiples of pi}}

\begin{corollary}\label{cor:algebraic-pi}
For every nonzero real algebraic number $\lambda$,
\[
 \boxed{
 \mu(\lambda\pi)
 \le[\QQ(i/\lambda):\QQ]
 \le2[\QQ(\lambda):\QQ].}
\]
If $\lambda^2\in\QQ^\times$, then $\mu(\lambda\pi)=2$.  For
example,
\[
 \mu(2^{1/3}\pi)\le6.
\]
\end{corollary}

\begin{proof}
Use the certificate
\[
 e^{(i/\lambda)(\lambda\pi)}=-1.
\]
Complex conjugation fixes $\lambda$ and sends $i/\lambda$ to its
negative, so the minimal field has two compatible embeddings.  The
first bound follows from \cref{thm:main}; the second follows from
$\QQ(i/\lambda)\subseteq\QQ(i,\lambda)$.  If $\lambda^2$ is rational,
then $\QQ(i/\lambda)$ is quadratic and both embeddings are compatible.
\end{proof}

\subsection{Arguments of algebraic numbers}

\begin{corollary}\label{cor:algebraic-arguments}
Let $z\in\Qbar^\times$, and let $\vartheta\in\RR^\times$ be an argument
of $z$.  Put
\[
 \beta=\frac z{\overline z}.
\]
Then
\[
 \boxed{
 \mu(\vartheta)
 \le
 \min_{N\ge1}
 [\QQ(i,\beta^N):\QQ].}
\]
In particular, for every nonzero real algebraic $\alpha$,
\[
 \boxed{
 \mu(\arctan\alpha)
 \le2[\QQ(\alpha):\QQ].}
\]
For rational $\alpha\ne0$, equality $\mu(\arctan\alpha)=2$ holds.
\end{corollary}

\begin{proof}
Since $e^{2i\vartheta}=\beta$, complex conjugation sends the coefficient
$2i$ to $-2i$ and $\beta$ to $\beta^{-1}$.  The powered certificate
therefore has two compatible embeddings in
$\QQ(i,\beta^N)$, so \cref{thm:main} gives the first bound.  For
$\vartheta=\arctan\alpha$, take
$z=1+i\alpha$ and use
$\QQ(i,\beta)\subseteq\QQ(i,\alpha)$.  The rational case has field
$\QQ(i)$ and complexity one.
\end{proof}

\subsection{Rational projective invariance}

\begin{proposition}[$\PGL_2(\QQ)$-invariance]\label{prop:pgl-invariance}
Let $x$ be irrational and let
\[
 g(X)=\frac{AX+B}{CX+D}\in\PGL_2(\QQ)
\]
be defined at $x$.  Then
\[
 \boxed{\mu(g(x))=\mu(x).}
\]
\end{proposition}

\begin{proof}
After multiplying the coefficients by a common denominator, assume
$A,B,C,D\in\ZZ$.  For a rational $p/q$ tending to $x$,
\[
 g(x)-g(p/q)
 =\frac{(AD-BC)(x-p/q)}
 {(Cx+D)(Cp/q+D)}.
\]
The two denominator factors stay bounded above and away from zero.  The
rational number
\[
 g(p/q)=\frac{Ap+Bq}{Cp+Dq}
\]
has reduced denominator $O(q)$.  Hence approximations to $x$ of every
exponent $\nu$ yield approximations to $g(x)$ of every smaller exponent.
Thus $\mu(g(x))\ge\mu(x)$.  Applying the same argument to $g^{-1}$ gives
the reverse inequality.
\end{proof}

Consequently every exact-two family above produces an entire rational
Möbius orbit of exact-two constants, including reciprocals and rational
affine transforms.

\section{Conjectures and further directions}\label{sec:conjectures}

Rigidity shows that the scalar product-formula method is now essentially
classified.  When $\kappa_{\exp}(\xi)>1$, no hidden collection of three
or more compatible scalar fibres can improve the bound.  Moreover,
\cref{prop:soft-barrier} shows that even the optimal aggregate use of the
approximation and collision energies remains subject to the same
threshold.  Any general passage from $2\kappa_{\exp}(\xi)$ to $2$ must
therefore use a genuinely non-scalar arithmetic or incidence input.

\begin{conjecture}[Algebraic exponential-fibre conjecture]
\label{conj:full-fibre}
If
\[
 \xi\in\RR^\times,
 \qquad
 a,b\in\Qbar^\times,
 \qquad
 e^{a\xi}=b,
\]
then
\[
 \boxed{\mu(\xi)=2.}
\]
\end{conjecture}

The first benchmark beyond the scalar theory is
\[
 \mu(\log\alpha)=2
 \qquad
 (\alpha\in\Qbar_{>0}\setminus\{1\}),
\]
for which \cref{thm:alglog-complexity} currently gives the optimal
scalar bound $2d_{\mathrm{pow}}(\alpha)$.

\begin{conjecture}[Algebraic linear forms in logarithms]
\label{conj:linear-forms}
Let $\alpha_1,\ldots,\alpha_s\in\Qbar^\times$, choose logarithm
branches, and let $\beta_1,\ldots,\beta_s\in\Qbar$.  If
\[
 \Lambda=\sum_{j=1}^s\beta_j\log\alpha_j
\]
is real and nonzero, then
\[
 \boxed{\mu(\Lambda)=2.}
\]
\end{conjecture}

Baker's theorem implies that such a nonzero $\Lambda$ is transcendental
\cite[Chapter~2]{Baker1975}; the conjecture asks for the optimal rational
approximation exponent.  It contains \cref{conj:full-fibre} as the
one-logarithm case.

\subsection{Full-Pl\"ucker arithmetic}

A plausible route to \cref{conj:full-fibre} is to retain the entire
vector of maximal minors of the interpolation matrix rather than one
scalar minor.  The scalar norm argument records only one Pl\"ucker
coordinate and is forced to pay the integer complexity
$\kappa_{\exp}$.  In contrast, \cref{cor:plucker-analytic} proves that the
analytic smallness estimate already holds uniformly for every maximal
minor and hence for every fixed norm of the full Pl\"ucker vector.  The
remaining problem is arithmetic: one needs an adelic projective-height
principle that converts simultaneous smallness of all maximal minors at
the compatible places into an exceptional functional relation without
paying separately for every incompatible scalar coordinate.

An adelic full-Pl\"ucker argument, potentially combined with the
Subspace Theorem \cite{Schmidt1980}, could produce such a relation.  The
moving-centre interpolation theorem is designed to exclude precisely
these functional dependencies.  The mixed-determinant perturbation
estimates of Litsey provide related evidence that families of minors can
be controlled stably under structured perturbations
\cite[Sections~5.2--5.6]{Litsey2026}.  Making the required projective
product-formula mechanism quantitative remains open.

\subsection{Near-compatible powers}

Let $L$ be a normal closure of $\QQ(a,b)$ and, for
$\tau:L\hookrightarrow\CC$, put
\[
 \delta_\tau=\frac{e^{\tau(a)\xi}}{\tau(b)}.
\]
The torsion-defect criterion \cref{cor:torsion-new} treats the case in
which every $\delta_\tau$ is a root of unity.

\begin{conjecture}[Unitary-defect exact-two principle]
\label{conj:unitary-defect}
If $|\delta_\tau|=1$ for every embedding $\tau:L\hookrightarrow\CC$,
then
\[
 \boxed{\mu(\xi)=2.}
\]
\end{conjecture}

The root-of-unity case follows from power saturation.  In the general
unitary case, simultaneous recurrence can make finitely many powers
$\delta_\tau^N$ close to $1$, but the determinant also amplifies the
remaining mismatch through the parameters $K$ and $H$.  A proof would
therefore require a quantitative near-compatible-fibre estimate together
with a coupled choice of the power $N$ and the interpolation scale.  The
fixed-certificate argument of the present paper does not provide such a
coupling.

\subsection{Effectivity of the geometric input}

\begin{conjecture}[Effective special-centre interpolation]
\label{conj:effective-interpolation}
For algebraic $a,b$, rational $x_1,\ldots,x_m$, and centres
\[
 y_j=b^j,
 \qquad
 c_{ji}=ja x_i,
 \qquad 0\le j<K,
\]
the eventual surjectivity threshold in \cref{thm:MI} is bounded
effectively in terms of the weights, $m$, $K$, the algebraic degrees, and
the logarithmic Weil heights of $a,b,x_1,\ldots,x_m$.
\end{conjecture}

Such a result would make the geometric part of the determinant argument
quantitative and would be an essential input for near-compatible powers
or explicit finite-place estimates.  It would also isolate whether the
remaining ineffectivity is intrinsic to the selection of rational
approximants or only to the present use of Serre vanishing.

\section*{Acknowledgments and provenance}

We thank Eric Naslund for publicly drawing attention, in posts on X on
9 and 12 October 2026, to the quadratic
$\arctanh(t\sqrt D)/\sqrt D$ family; the later post linked a public
Claude discussion \cite{Naslund2026}.  We cite those posts as prior
public identification of the family.  The converse rigidity, intrinsic
power-saturation, and higher-degree bounds are contributions of the
present paper.

The geometric interpolation theorem and the analytic row-translation
and Taylor-collision estimates used here are those of Liu--Jiang--Zhang
\cite{LiuJiangZhang2026}, building on the interpolation-determinant
method of Laurent \cite{Laurent1991,Laurent2000} and the proof for
$\pi$ in \cite{OpenAI2026}.  Brown independently obtained the
rational-logarithm theorem and a general-fibre interpolation theorem
under a global fibre-count condition, together with a Lean~4
formalization of the rational theorem, and recorded the real-algebraic
degree bound \cite{Brown2026}.  Litsey's mixed-determinant perturbation
argument for Catalan's constant provides related contemporary context
for continuous defect optimization \cite{Litsey2026}.  The present
paper combines the Liu--Jiang--Zhang approximation and collision
penalties by the soft convex optimization of
\cref{lem:collision-energy}; the new arithmetic determinant step is the
simultaneous number-field denominator clearing and product-formula
comparison across compatible and incompatible conjugate fibres.  The
rigidity, intrinsic power-saturation, and torus classification are
obtained by combining that comparison with Gelfond--Schneider.

\renewcommand*{\bibfont}{\small}
\printbibliography

\end{document}
